
\chapter{Modelado de amenazas mediante STRIDE}

\section{Metodología STRIDE}

STRIDE es una metodología de modelado de amenazas propuesta por Microsoft
\cite{stride}. Parte de un diagrama de flujo de datos (DFD) del sistema y
revisa cada elemento del diagrama (actores externos, procesos, almacenes de
datos y flujos) frente a seis categorías de amenaza. Cada categoría
corresponde a la violación de una propiedad de seguridad: autenticidad,
integridad, no repudio, confidencialidad, disponibilidad y autorización.

En MoonPhases se aplicó en cuatro pasos:

\begin{enumerate}
    \item \textbf{Verificación de la arquitectura} contra el repositorio
    (\texttt{compose.yaml}, \texttt{nginx.conf}, \texttt{server.js},
    \texttt{db.js} y los \texttt{Dockerfile}), para no modelar componentes
    que no existen.
    \item \textbf{Construcción del DFD} con actores, procesos, almacenes,
    flujos y límites de confianza (\textit{Trust Boundaries}).
    \item \textbf{Evaluación por elemento:} para cada componente o flujo
    se consideró qué categorías STRIDE le son aplicables. No se asumió que
    toda la aplicación sea vulnerable a las seis categorías.
    \item \textbf{Clasificación de cada amenaza} por su estado
    (\textit{confirmada}, \textit{candidata} o \textit{mitigada}) y por su
    nivel de riesgo, calculado como impacto por probabilidad.
\end{enumerate}

Las seis categorías consideradas son:

\begin{table}[H]
\centering
\caption{Categorías de amenazas STRIDE.}
\label{tab:stride}

\begin{tabularx}{\textwidth}{l l X}
\toprule
\textbf{Categoría} &
\textbf{Abreviatura} &
\textbf{Descripción} \\
\midrule
Suplantación de identidad & S &
Suplantación de una entidad legítima. \\

Manipulación de datos & T &
Modificación no autorizada de información. \\

Repudio & R &
Negación de acciones realizadas. \\

Divulgación de información & I &
Acceso no autorizado a información. \\

Denegación de servicio & D &
Interrupción de la disponibilidad. \\

Elevación de privilegios & E &
Obtención de permisos superiores. \\
\bottomrule
\end{tabularx}
\end{table}

\subsection{Criterio de riesgo}

El nivel de riesgo no se asignó de forma arbitraria. Se calculó como
\textbf{riesgo = impacto $\times$ probabilidad}, con ambos factores en una
escala de 1 a 3:

\begin{table}[H]
\centering
\caption{Escalas utilizadas para calcular el riesgo.}
\label{tab:criterio-riesgo}

\begin{tabularx}{\textwidth}{c X X}
\toprule
\textbf{Valor} &
\textbf{Impacto} &
\textbf{Probabilidad} \\
\midrule
1 &
Facilita otros ataques; no hay pérdida directa. &
Requiere acceso previo (host, cuenta o compromiso del backend). \\

2 &
Degradación del servicio o pérdida parcial de datos. &
Requiere una condición adicional (estar en la ruta de red, un sitio
visitado por la víctima, etc.). \\

3 &
Alteración o pérdida total de los datos, o control del componente. &
Explotable hoy desde la red, sin precondiciones. \\
\bottomrule
\end{tabularx}
\end{table}

Un puntaje de 1 a 2 se considera \textbf{bajo}, de 3 a 4 \textbf{medio} y de
6 a 9 \textbf{alto}. Las amenazas se clasifican además por su estado:
\textit{confirmada} cuando hay evidencia directa en el código o en la
configuración, \textit{candidata} cuando depende de una condición que no se
ha verificado, y \textit{mitigada} cuando existe un control con riesgo
residual.

\section{Diagrama de modelado de amenazas}

La figura~\ref{fig:stride} presenta el diagrama de flujo de datos con el
análisis STRIDE, elaborado en draw.io. Muestra únicamente la arquitectura
actual del sistema, verificada contra el repositorio. Las mitigaciones
propuestas se describen en el capítulo de hallazgos y no aparecen dibujadas.
El archivo editable (\texttt{.drawio}) y la matriz completa de amenazas se
encuentran en \texttt{docs/stride/}.

\begin{landscape}
\evidencefigure
{figures/stride.pdf}
{Diagrama de flujo de datos de MoonPhases con el análisis STRIDE.}
{fig:stride}
\end{landscape}

\subsection{Actores y componentes}

\begin{itemize}
    \item \textbf{Actores externos:} el usuario legítimo, que accede desde un
    navegador; el atacante externo, que puede intentar explotar la
    aplicación por el camino normal o directamente; y el desarrollador, que
    administra el repositorio de GitHub.
    \item \textbf{Procesos:} el navegador con la aplicación de Three.js (P0);
    Nginx sirviendo el frontend estático (P1); el mismo Nginx como proxy
    inverso de \texttt{/api/} (P2); la API REST de Express (P3); el gestor
    de conexiones \texttt{pg.Pool} (P4); el servidor PostgreSQL (P5); y el
    repositorio de GitHub (P6).
    \item \textbf{Almacenes de datos y activos:} la tabla
    \texttt{moon\_notes} (DS1), el volumen persistente
    \texttt{postgres-data} (DS2), el archivo \texttt{pg\_password.txt} con el
    secreto montado en los contenedores (DS3) y el código fuente (DS4).
\end{itemize}

Durante la verificación se corrigió un dato: el frontend publica el puerto
\textbf{8081}, no el 8080. Lo confirman \texttt{listen 8081} en
\texttt{nginx.conf} y \texttt{8081:8081} en \texttt{compose.yaml}. El backend
escucha en el 3001 y PostgreSQL en el 5432, solo dentro de la red privada.
No se dibujaron servicios que el proyecto no implementa.

\subsection{Flujos de datos}

\begin{table}[H]
\centering
\caption{Flujos de datos representados en el diagrama.}
\label{tab:flujos}

\begin{tabularx}{\textwidth}{l X}
\toprule
\textbf{Flujo} &
\textbf{Descripción} \\
\midrule
F1, F2 &
El navegador solicita la aplicación a Nginx por HTTP (puerto 8081) y
recibe la aplicación estática y las respuestas. \\

F3 &
El navegador envía las operaciones sobre notas al proxy mediante
\texttt{/api/}. \\

F4 &
Nginx reenvía la solicitud a Express con \texttt{proxy\_pass}. \\

F5, F6 &
Express consulta PostgreSQL con SQL parametrizado, sin TLS, y recibe los
resultados a través del pool de conexiones. \\

F7 &
Docker Compose entrega la contraseña de PostgreSQL al backend y a la base de
datos como archivo montado. \\

F8 &
El desarrollador sube código al repositorio de GitHub. \\

F9 &
El atacante intenta llegar a la API por el camino normal (Nginx). \\

F10 &
El atacante accede directamente al puerto 3001, publicado en el host por
\texttt{compose.yaml}. Se dibuja en rojo porque evita el control de Nginx.
Es un flujo confirmado en la configuración, no una suposición. \\
\bottomrule
\end{tabularx}
\end{table}

\subsection{Límites de confianza}

\begin{table}[H]
\centering
\caption{Trust Boundaries del modelo.}
\label{tab:tb}

\begin{tabularx}{\textwidth}{l X l}
\toprule
\textbf{ID} &
\textbf{Qué separa} &
\textbf{Tipo} \\
\midrule
TB1 & Internet y navegador frente al servidor de la aplicación. &
Red \\

TB2 & Red Docker pública, donde están Nginx y Express. &
Red \\

TB3 & Red Docker privada (\texttt{internal: true}), donde Express se
comunica con PostgreSQL. &
Red \\

TB4 & Entorno de desarrollo y repositorio de GitHub. &
Lógica \\

TB5 & Acceso a los secretos y al almacenamiento persistente del host. &
Lógica \\
\bottomrule
\end{tabularx}
\end{table}

TB1 a TB3 corresponden a fronteras de red reales. TB4 y TB5 son zonas
lógicas: no son redes, pero delimitan quién puede acceder al repositorio y a
los archivos del host. El backend pertenece a la red pública y también está
conectado a la privada; en el diagrama se ubica en TB2 y el flujo F5 cruza
hacia TB3.

\subsection{Ubicación de las amenazas}

Cada amenaza se dibuja junto al componente o flujo que afecta. Se agrupan en
Express (P3) las que dependen de su código (autenticación, validación, CORS,
errores y registro), y en Nginx (P1 y P2) las del canal y los límites de
tasa. La amenaza S2 se dibuja sobre el flujo F10 porque la amenaza es el
camino mismo. Las demás se colocan en el almacén o proceso afectado. El color
del círculo indica el nivel de riesgo y el contorno discontinuo indica que la
amenaza es candidata o está mitigada, no confirmada.

\section{Principales amenazas identificadas}

El modelado identificó 16 amenazas. Las cinco más relevantes (seis
identificadores, porque la primera agrupa S1 y T1) se presentan a
continuación, ordenadas por riesgo y por el respaldo que tienen en el código
y en las herramientas. Para cada una se indica si es una vulnerabilidad
\textbf{confirmada por herramientas} o solo una amenaza deducida del código y
de la configuración. El resto se resume al final de esta sección.

% ------------------------------------------------------------
% AMENAZA 1
% ------------------------------------------------------------

\subsection{Amenaza 1 (S1 y T1): API sin autenticación ni autorización}

\begin{itemize}
    \item \textbf{Categoría STRIDE:} Suplantación (S) y Manipulación (T).
    \item \textbf{Componente afectado:} API REST de Express (P3) y tabla
    \texttt{moon\_notes} (DS1).
    \item \textbf{Descripción:} ninguna ruta verifica la identidad de quien
    llama. Cualquier cliente que alcance la API puede listar, crear,
    sobrescribir (\texttt{POST} hace un \textit{upsert}) y borrar la nota de
    cualquier fecha. El propio esquema de la base de datos declara ``sin
    autenticación de usuario''.
    \item \textbf{Impacto:} alteración o pérdida de los datos del diario.
    Cualquier visitante actúa con el máximo privilegio disponible.
    \item \textbf{Estado y evidencia:} confirmada por revisión del código
    (\texttt{server.js}, \texttt{schema.sql}). Riesgo \textbf{alto (9)}:
    impacto 3 por probabilidad 3. ZAP, además, accede a \texttt{/api/notes}
    sin credenciales, lo que es coherente con la amenaza, aunque el escaneo
    pasivo no intenta modificar datos.
    \item \textbf{Mitigación:} añadir autenticación (sesión o JWT) y
    autorización por recurso. Como mínimo, proteger \texttt{POST} y
    \texttt{DELETE} con un token.
\end{itemize}

% ------------------------------------------------------------
% AMENAZA 2
% ------------------------------------------------------------

\subsection{Amenaza 2 (D1): ausencia de límites de tasa y de tamaño}

\begin{itemize}
    \item \textbf{Categoría STRIDE:} Denegación de servicio (D).
    \item \textbf{Componente afectado:} proxy inverso de Nginx (P2) y API
    de Express (P3).
    \item \textbf{Descripción:} no hay \texttt{limit\_req} en Nginx ni límite
    de longitud para el campo \texttt{content}. Solo rigen los valores por
    defecto de Nginx (1~MB por cuerpo) y de \texttt{express.json()}
    (100~KB). Un cliente puede enviar ráfagas de peticiones o notas grandes
    en bucle.
    \item \textbf{Impacto:} saturación del backend, del pool de conexiones o
    del volumen de PostgreSQL.
    \item \textbf{Estado y evidencia:} confirmada por revisión de
    \texttt{nginx.conf} y \texttt{server.js}. Ninguna de las herramientas
    implementadas la detecta, porque no hay una prueba de carga contra el
    servicio. Riesgo \textbf{alto (6)}: impacto 2 por probabilidad 3.
    \item \textbf{Mitigación:} \texttt{limit\_req} y
    \texttt{client\_max\_body\_size} en Nginx, y validación de la longitud de
    \texttt{content} en Express.
\end{itemize}

% ------------------------------------------------------------
% AMENAZA 3
% ------------------------------------------------------------

\subsection{Amenaza 3 (T3): CORS abierto a cualquier origen}

\begin{itemize}
    \item \textbf{Categoría STRIDE:} Manipulación (T).
    \item \textbf{Componente afectado:} API de Express (P3).
    \item \textbf{Descripción:} \texttt{app.use(cors())} sin opciones
    responde \texttt{Access-Control-Allow-Origin: *}. Cualquier sitio web
    que visite el usuario puede invocar la API desde su navegador.
    \item \textbf{Impacto:} lectura y escritura de notas desde páginas de
    terceros, favorecida por la falta de autenticación (amenaza 1).
    \item \textbf{Estado y evidencia:} \textbf{confirmada por herramienta}.
    OWASP ZAP la reporta como \textit{Cross-Domain Misconfiguration}, con la
    cabecera \texttt{Access-Control-Allow-Origin: *} en \texttt{/api/}.
    Riesgo \textbf{medio (4)}.
    \item \textbf{Mitigación:} \texttt{cors(\{ origin: <origen del
    frontend> \})}.
\end{itemize}

% ------------------------------------------------------------
% AMENAZA 4
% ------------------------------------------------------------

\subsection{Amenaza 4 (I1): divulgación de información en errores}

\begin{itemize}
    \item \textbf{Categoría STRIDE:} Divulgación de información (I).
    \item \textbf{Componente afectado:} API de Express (P3), respuestas
    HTTP 500.
    \item \textbf{Descripción:} los cuatro bloques \texttt{catch} devuelven
    \texttt{err.message} al cliente. Un error de PostgreSQL, por ejemplo con
    una fecha mal formada, revela detalles del motor.
    \item \textbf{Impacto:} facilita el reconocimiento de la aplicación y
    ataques posteriores. No da acceso directo a los datos.
    \item \textbf{Estado y evidencia:} confirmada por revisión del código y
    \textbf{respaldada por herramienta}: ZAP detectó una respuesta
    \texttt{500 Internal Server Error} en \texttt{/api/notes}
    (\textit{Application Error Disclosure}) y la cabecera
    \texttt{X-Powered-By: Express}, que revela la tecnología del backend.
    Riesgo \textbf{medio (3)}: impacto 1 por probabilidad 3.
    \item \textbf{Mitigación:} registrar el error completo en el servidor,
    responder con un mensaje genérico y desactivar \texttt{X-Powered-By}.
\end{itemize}

% ------------------------------------------------------------
% AMENAZA 5
% ------------------------------------------------------------

\subsection{Amenaza 5 (T4): canal sin TLS y sin cabeceras de seguridad}

\begin{itemize}
    \item \textbf{Categoría STRIDE:} Manipulación (T).
    \item \textbf{Componente afectado:} Nginx (P1) y los flujos F1 y F2.
    \item \textbf{Descripción:} el tráfico viaja por HTTP sin TLS y Nginx no
    envía cabeceras de seguridad. Las respuestas pueden manipularse en
    tránsito y la aplicación puede incrustarse en un \texttt{iframe}
    (\textit{clickjacking}).
    \item \textbf{Impacto:} alteración del contenido que recibe el usuario y
    ausencia de políticas que limiten la ejecución de scripts.
    \item \textbf{Estado y evidencia:} \textbf{confirmada por herramienta}.
    ZAP reporta la ausencia de \textit{Content Security Policy},
    \textit{Anti-clickjacking Header} y \textit{X-Content-Type-Options}.
    Además, marca la versión de Nginx como vulnerable
    (\textit{nginx Version with Known Vulnerabilities}). Riesgo \textbf{medio
    (4)}.
    \item \textbf{Mitigación:} TLS en Nginx, cabeceras CSP,
    \texttt{X-Frame-Options}, \texttt{X-Content-Type-Options},
    \texttt{Referrer-Policy} y HSTS, y actualizar la imagen de Nginx.
\end{itemize}

\subsection{Otras amenazas del modelo}

Las diez amenazas restantes tienen un riesgo menor o dependen de
condiciones que no se han verificado. Se resumen en la
tabla~\ref{tab:otras-amenazas}.

\begin{table}[H]
\centering
\caption{Otras amenazas identificadas en el modelo.}
\label{tab:otras-amenazas}

\begin{tabularx}{\textwidth}{l l X l l}
\toprule
\textbf{ID} &
\textbf{Categoría} &
\textbf{Amenaza} &
\textbf{Estado} &
\textbf{Riesgo} \\
\midrule
S2 & S & Puerto 3001 publicado: se accede al backend sin pasar por Nginx. &
Confirmada & Medio (3) \\

S3 & S & Cuenta del desarrollador comprometida; \texttt{main} exige 1
revisión por PR. & Mitigada & Medio (3) \\

T2 & T & \texttt{:date} y \texttt{content} sin validar; el SQL
parametrizado evita la inyección. & Confirmada & Medio (4) \\

T5 & T & Imágenes base por tag, no por digest (cadena de suministro). &
Confirmada & Medio (3) \\

R1 & R & Escrituras y borrados sin registro de auditoría. &
Confirmada & Medio (4) \\

I2 & I & Contraseña en archivo plano en el host (Docker secret). &
Mitigada & Medio (3) \\

I3 & I & Conexión a PostgreSQL sin TLS, en red interna. &
Candidata & Bajo (2) \\

I4 & I & Volumen \texttt{postgres-data} sin cifrado en reposo. &
Candidata & Bajo (2) \\

D2 & D & \texttt{pg.Pool} con valores por defecto, sin timeouts. &
Candidata & Medio (4) \\

E1 & E & El rol \texttt{moonuser} es superusuario en la imagen oficial de
PostgreSQL. & Candidata & Medio (3) \\
\bottomrule
\end{tabularx}
\end{table}

En la API no hay escalamiento de privilegios porque no existen roles:
quien llega a ella ya tiene el máximo permiso, y eso se recoge en la
amenaza 1 (S1 y T1). Por eso la única amenaza de elevación de privilegios
es E1, condicionada a que el backend ya esté comprometido.

Existen también controles que reducen el riesgo y deben conservarse:
consultas SQL parametrizadas, PostgreSQL sin publicar en el host y en una
red interna, procesos sin permisos de administrador, contraseña como Docker
secret fuera de Git, dependencias fijadas por lockfile y protección de la
rama \texttt{main}.

\section{Resumen de amenazas}

\begin{table}[H]
\centering
\caption{Resumen de las principales amenazas STRIDE.}
\label{tab:amenazas-principales}

\begin{tabularx}{\textwidth}{c l X l}
\toprule
\textbf{ID} &
\textbf{Categoría} &
\textbf{Amenaza} &
\textbf{Prioridad} \\
\midrule
S1, T1 & Suplantación y Manipulación &
API sin autenticación ni autorización: cualquiera lee, crea, modifica y
borra notas. (Confirmada en código) & Alta \\

D1 & Denegación de servicio &
Sin límites de tasa ni de tamaño de contenido. (Confirmada en código) &
Alta \\

T3 & Manipulación &
CORS abierto a cualquier origen. (Confirmada por ZAP) & Media \\

I1 & Divulgación de información &
Errores internos y tecnología expuestos en las respuestas HTTP.
(Confirmada en código y respaldada por ZAP) & Media \\

T4 & Manipulación &
HTTP sin TLS y sin cabeceras de seguridad. (Confirmada por ZAP) & Media \\
\bottomrule
\end{tabularx}
\end{table}
